%% ma: L12-B08-0104
%% lop: 12
%% chuong: 2
%% bai: 8
%% dang: 12.8.7
%% loai: TN4
%% muc: NB
%% dap_an: "A"
%% nguon: "CĐ2.3.pdf giáo viên gửi 27/10/2026 (M07, Lê Minh Khôi – Nguyễn Anh Khôi, Chương II Chủ đề 3 Biểu thức tọa độ của các phép toán vectơ), Câu 15"
%% kien_thuc: "Bài 8 (tọa độ trung điểm)"
%% trang_thai: chinh-thuc
\begin{ex}
Trong không gian $Oxyz$, cho hai điểm $D(-9;1;-12)$ và $E(-5;3;2)$. Tọa độ của điểm $I$ thỏa mãn $\overrightarrow{ID}+\overrightarrow{IE}=\vec 0$ là
\choice{\True $I(-7;2;-5)$}{$I(7;-2;5)$}{$I\left(-\dfrac72;1;-\dfrac52\right)$}{$I\left(-\dfrac{14}3;\dfrac43;-\dfrac{10}3\right)$}
\loigiai{$I$ là trung điểm $DE$: $I(-7;2;-5)$.}
\end{ex}
